\section{Conclusion}
\label{sec:conclusion}

% ---- 脉络（写完删除）--------------------------------------------------------
% 一段：重申"生成形状 vs 控制形状"的论点 -> 三个组件 -> 主结果 -> 对 WAM 设计的启示。
% Limitations（可并入本节或单独成段）：
%   - 结论依赖运行点（future 帧近似噪声占位是 tau_v = 1 时的事实）；
%   - 分配在套件内相对稳定，价值在"放在哪里"而非"逐步追踪场景"；
%   - 蒸馏目标的最优解即 teacher，"不可区分"是该目标下的上限；
%   - 时间轴仍有空间（真实未来运动 oracle 离线更优，但推理时不可得）；
%   - 主要在单一 backbone 上验证。
% ---------------------------------------------------------------------------

\TODO{Conclusion and limitations.}
